% =============================================================================
%  Sección 2: Desarrollo
%  Sustituye el contenido de \lipsum con tu propio texto.
%  Agrega o elimina subsecciones según la estructura de tu argumento.
% =============================================================================

\section{Marco Conceptual y Desarrollo Argumentativo}

\subsection{Fundamentos teóricos}

\lipsum[5]

\citaclave{
  ``La ciudad compacta no es una utopía formal, sino una decisión política sobre cómo se distribuyen
  los recursos del suelo urbano entre quienes lo necesitan y quienes lo acaparan.''
  \citep{lore2024densidad}
}

\lipsum[6]

\subsection{Modelos de referencia internacional}

\lipsum[7-8]

% ── Tabla de ejemplo ─────────────────────────────────────────────────────────
\begin{table}[ht]
\centering
\caption{Comparativa de modelos de densificación urbana.}
\label{tab:modelos}
\begin{tabularx}{\textwidth}{@{} l X X X @{}}
\toprule
\textbf{Modelo}    & \textbf{Principio rector} & \textbf{Instrumento clave} & \textbf{Caso referente} \\
\midrule
\rowcolor{ColorFondoTabla}
TOD Nórdico        & Proximidad al transporte  & Plan maestro multiescala  & Copenhague, Estocolmo \\
Barrio Consolidado & Rehabilitación incremental& Participación comunitaria & Medellín, Bogotá \\
\rowcolor{ColorFondoTabla}
Corredor Urbano    & Mixticidad y conectividad & Corredores de inversión   & Ciudad de México \\
\bottomrule
\end{tabularx}
\fuentefigura{Elaboración propia con base en \citet{lore2024densidad} y \citet{lore2024zmvm}.}
\end{table}

\subsection{Aplicación al contexto mexicano}

\lipsum[9-10]

% ── Figura de ejemplo ────────────────────────────────────────────────────────
% Para agregar una figura real, descomenta y ajusta:
%
% \begin{figure}[ht]
%   \centering
%   \includegraphics[width=0.8\textwidth]{img/nombre_figura.png}
%   \caption{Descripción de la figura.}
%   \fuentefigura{Fuente de la imagen.}
%   \label{fig:nombre_label}
% \end{figure}

\lipsum[11]
